%En este Trabajo de Fin de Máster estudiaremos el problema de los factores compartidos en claves públicas RSA extraídas de \textit{Certificate Transparency logs}. Esta vulnerabilidad puede aparecer cuando la generación de claves no se realiza con suficiente aleatoriedad, provocando que distintos módulos RSA compartan un mismo factor primo y puedan factorizarse de forma eficiente.


\cambio{El algoritmo RSA es uno de los esquemas de cifrado asimétrico más populares hasta la fecha, del que se han estudiado múltiples vulnerabilidades.  Una de ellas puede producirse cuando la generación de claves RSA no se realiza con suficiente aleatoriedad, provocando que distintos módulos RSA compartan un mismo factor primo y puedan factorizarse de forma eficiente. }\commA{Esto solo lo he cambiado de orden, es mejor empezar con algo general e ir luego a lo particular}
\\
\\
%El trabajo se apoyará, por un lado, 
\cambio{La explotación de esta vulnerabilidad fue estudiada} en la tesis \textit{Finding shared RSA factors in the Certificate Transparency logs}, que abordó este problema mediante un enfoque clásico de \textit{batch GCD} \commA{Estaría bien que comentaras en pocas palabras este enfoque para que sea el resumen autocontenido }, y, por otro, Pelofske, propone el algoritmo \textit{Binary Tree Batch GCD} como una alternativa más eficiente. A partir de estas referencias, implementaremos dicho algoritmo con el objetivo de estudiar su utilidad práctica en este contexto.
\\
\\
\cambio{En este trabajo} hemos desarrollado una implementación en C++17 con GMP y OpenMP, que ya ha sido validada en una primera fase sobre datos sintéticos generados \cambio{expresamente para demostrar la  capacidad del algoritmo}.\commA{No me convence el decir solo de forma controlada} Esta etapa nos ha permitido comprobar que el algoritmo detecta correctamente factores compartidos cuando estos existen y observar su comportamiento \cambio{en términos de rendimiento computacional} \commA{no me convence decir solo que rendimiento temporal, tal vez puedas redactarlo de otra forma?} al aumentar el tamaño del conjunto de entrada.
%de forma controlada
%\\
%\\
Posteriormente, hemos aplicado la misma implementación sobre una primera muestra real de módulos RSA-2048 extraídos de \textit{CT logs}. Aunque esta prueba todavía es inicial y no equivale a un análisis masivo de todos los logs, sí constituye una validación útil del procedimiento sobre datos reales.
\\
\\
En una fase posterior del trabajo ampliaremos este análisis sobre un conjunto mucho mayor de \textit{CT logs}, con el fin de evaluar mejor la escalabilidad del método y su utilidad para auditorías criptográficas a gran escala. En conjunto, el TFM combinará estudio de trabajos previos, desarrollo metodológico, implementación experimental y análisis aplicado sobre un problema real de ciberseguridad. \commA{Aqui tendrías que explicar si es trabajo futuro, o es algo que vas a hacer en la memoria }


